\section{Introduction}

Capacity-constrained clustering gives BISECT a construction family whose
primary abstraction is assignment to population-capacity clusters. This makes
capacity status, repair, and infeasibility visible rather than hidden inside a
single final assignment.

The central design choice is to make failure part of the artifact contract.
Many construction algorithms only expose a final plan or an error. T.15 exposes
the intermediate accountability surface: which deterministic seeds were chosen,
which assignment rule was used, whether capacity was satisfied, whether repair
was needed, and whether the exported plan is backed by RPLAN and RCTX sidecars.
That makes the paper useful even before broad empirical comparison, because it
defines how a clustered construction should explain itself.
